%% SCRAP: experiments/02-experiments/physics-optimization/proposals
%% SOURCE: docs/working/experiments/02-experiments/physics-optimization/proposals.md
%% STATUS: CURRENT
%% FITS: experiments/ch-physics-opt, vol1-vm-physics/ch-loops
%% EDITORIAL: lifted — prose rewritten to press voice

\section{Physics-Driven Optimisation Proposals}
\label{sec:physics-opt-proposals}

\subsection{Strategic Direction}

The hot-words cache experiment (1.78$\times$, 95\% CI: [1.75$\times$,
1.81$\times$]) validates the physics-driven optimisation methodology:
collect execution-frequency metrics at runtime, make automatic promotion
decisions, and measure impact with statistical rigour. Nine additional
opportunities build on the same foundation.

JIT compilation is explicitly excluded from consideration. JIT requires
runtime code generation, violates L4Re microkernel policy, introduces
floating-point overhead, and defeats formal verification. Physics-driven
optimisation achieves measurable gains within StarForth's verification and
portability constraints.

\subsection{Opportunity Catalogue}

\paragraph{Opportunity~1 — Return Stack Prediction and Colon Word Inlining.}
Frequently called colon definitions with small bodies (fewer than 256~bytes)
are candidates for compile-time inlining when \texttt{execution\_heat > 100}.
Inlining eliminates dictionary lookup and return-stack overhead.
Expected gain: 1.3--2.0$\times$ for call-heavy workloads. Verification:
static (inlining correctness is provable via Isabelle/HOL).

\paragraph{Opportunity~2 — Block I/O Prefetching.}
A Markov-style transition matrix tracks which block LBN follows which.
When a block is accessed and the next-block heat exceeds a threshold, that
block is pre-fetched into the buffer cache. Expected gain: 1.5--3.0$\times$
for block-sequential access patterns. Implementation: medium complexity.

\paragraph{Opportunity~3 — Stack Operation Fusion.}
Co-execution heat tracks consecutive word pairs (e.g., \texttt{DUP DROP},
\texttt{SWAP ROT}, \texttt{OVER SWAP}). Pairs exceeding a heat threshold
at compile time are fused into dedicated primitives, eliminating two lookups
and one execution per pair. Expected gain: 1.2--1.5$\times$ for stack-heavy
programs. Verification: static (composition of verified primitives).

\paragraph{Opportunity~4 — Memory Allocation Pattern Prediction.}
Allocation frequency by size is tracked in a heat vector. Sizes exceeding
a threshold trigger pre-warmed pool maintenance. Repeating allocation
sequences are detected and cached. Expected gain: 1.3--1.8$\times$ for
allocation-heavy programs.

\paragraph{Opportunity~5 — Control Flow Branch Prediction.}
\texttt{IF}/\texttt{THEN}/\texttt{ELSE} outcomes are tracked per source
location. Branches with greater than 90\% one-sided bias are annotated
with a prediction hint that the inner interpreter can use to reorder code
or emit x86 branch-hint prefixes. Expected gain: 1.1--1.3$\times$
(architecture-dependent).

\paragraph{Opportunity~6 — Vocabulary Search Path Reordering.}
Per-vocabulary hit rates are tracked. The search order is sorted
dynamically by hit rate, placing the most productive vocabulary first.
Expected gain: 1.2--1.6$\times$ in multi-vocabulary programs.
Implementation: low complexity.

\paragraph{Opportunity~7 — String Operation Batching.}
Consecutive string-output operations (\texttt{."}, \texttt{EMIT},
\texttt{TYPE}) detected at compile time are fused into a single batch
output, reducing interpreter loop iterations. Expected gain:
1.1--1.4$\times$ for I/O-bound programs. Implementation: low complexity.

\paragraph{Opportunity~8 — Arithmetic Operation Reordering.}
For hot commutative operations at a given source location, operand sizes
are compared and the smaller operand is loaded first to improve prefetch
behaviour. Expected gain: 1.05--1.15$\times$ (architecture-dependent).

\paragraph{Opportunity~9 — Word Placement Optimisation.}
Co-execution heat between word pairs guides memory compaction: frequently
co-executed words are relocated adjacent to each other in the dictionary
to improve instruction-cache locality. Expected gain: 1.05--1.2$\times$.
Implementation: high complexity (requires GC integration).

\subsection{Implementation Roadmap}

\begin{center}
\begin{tabular}{lll}
\toprule
Phase & Opportunities & Rationale \\
\midrule
1 (complete) & Hot-words cache & Proven; production-ready \\
2 (recommended next) & \#3, \#6, \#1 & Highest ROI, lowest complexity \\
3 (future) & \#2, \#7, \#4 & Medium complexity \\
4 (advanced) & \#5, \#8, \#9 & CPU-specific or GC-dependent \\
\bottomrule
\end{tabular}
\end{center}

\subsection{Expected Cumulative Gains}

Conservative sequential composition of Phases~1--2:

\[
  1.78\times \;\times\; 1.2\;\times\; 1.2\;\times\; 1.3
  \;\approx\; 3.3\times \text{ cumulative speedup}
\]

All nine opportunities together: 5--8$\times$ total improvement is
plausible, subject to workload dependency and diminishing returns.

\subsection{Unified Metrics Infrastructure}

All nine opportunities require co-execution heat tracking and sequence
pattern buffers not yet present in the VM. A one-time extension to the
\texttt{PhysicsMetrics} per-word structure adds: a co-execution heat
vector, timing accumulator, per-word cache-line counters, and a small
recent-execution circular buffer. This single investment unlocks the
infrastructure required for Opportunities~1--9.

